\section{Minecraft Analysis}

\subsection{Alignment with Programming Concepts}

Minecraft offers an interactive digital environment that closely aligns with core programming and computational thinking concepts. 
Actions performed by players occur in a clear sequence, which reflects the execution of instructions in imperative programming. 
Repetitive gameplay activities such as mining blocks, farming resources, or crafting items correspond to loop structures, where the same operation is repeated until a specific condition is satisfied. 
Decision-making in response to in-game situations, including reacting to nightfall, selecting appropriate tools, or avoiding hostile entities, mirrors conditional logic used in programming. 
Managing inventory items is similar to dealing with variables that change over time, while recurring building patterns and construction methods resemble the use of procedures and modular design. 
In-game events such as mob spawning or environmental changes parallel event-driven programming models. 
These conceptual parallels have been identified and discussed in multiple studies examining the educational use of Minecraft in STEM and computational thinking contexts \cite{md_shamsudin_crafting_2024, nebel_mining_2016}.

\subsection{Empirical Evidence for Computational Thinking Development}

Empirical evidence supports the effectiveness of Minecraft for developing computational thinking when its use is structured and guided. 
A study focusing on Minecraft-based coding activities with middle school students reported statistically significant improvements in computational thinking dimensions such as sequencing, loops, conditionals, events, and parallelism following a structured instructional intervention \cite{tablatin_using_2023}. 
Systematic reviews of empirical studies from the past decade also find that Minecraft can help learners grasp basic programming concepts, provided that learning goals and instructional support are clearly defined \cite{nebel_mining_2016, md_shamsudin_crafting_2024}.

\subsection{Adoption in Formal Education}

Minecraft has seen substantial adoption in formal educational contexts, primarily through Minecraft Education Edition, which is designed specifically for classroom use. 
According to Microsoft Education, tens of millions of teachers and students across more than one hundred countries have access to licensed versions of the platform, indicating widespread institutional uptake rather than isolated experimentation \cite{noauthor_licensing_nodate}. 
A systematic literature review covering studies published between 2015 and 2024 confirms that Minecraft has been implemented across primary and secondary education, most frequently within STEM subjects, and that many educators continue to use it beyond initial trials \cite{md_shamsudin_crafting_2024}. 
Although precise global percentages of schools using Minecraft are not consistently reported, the scale of adoption and the volume of peer-reviewed research demonstrate that Minecraft is an established educational platform.

\subsection{Educational Effectiveness}

Research on the educational effectiveness of Minecraft consistently reports increased student engagement, motivation, and participation when the platform is integrated into structured learning activities. 
Studies grounded in constructivist and social constructivist learning theories highlight that Minecraft supports inquiry-based learning, problem solving, collaboration, and spatial reasoning \cite{short_teaching_2012, kotsopoulos_pedagogical_2017}. 
Experimental research further indicates that Minecraft-based interventions can lead to measurable improvements in spatial and cognitive skills, provided that instructional goals are explicit and teacher guidance is present \cite{ming_exploring_2025}.

\subsection{Limitations and Technical Barriers}

At the same time, the literature identifies significant limitations related to technical and programming complexity. 
Multiple studies report that educators often struggle to design effective Minecraft-based lessons due to limited programming knowledge and lack of confidence with the platform's technical features \cite{bower_improving_2017, nebel_mining_2016}. 
A recent systematic review emphasizes that insufficient teacher training and high technical barriers frequently result in Minecraft being used in a superficial manner, which reduces its potential to support deeper learning outcomes \cite{md_shamsudin_crafting_2024}. 
These findings indicate that Minecraft's educational benefits are strongly dependent on pedagogical structure and accessible tooling.

\subsection{Implications for DSL Design}

The implications of these findings are directly relevant to the development of a domain-specific language for children's education. 
The reviewed research demonstrates that Minecraft provides a motivating and conceptually appropriate environment for introducing computational thinking, yet its effectiveness is constrained by technical complexity and the programming expertise required from both teachers and learners \cite{md_shamsudin_crafting_2024}. 
A simplified domain-specific language can address these constraints by reducing syntactic and cognitive load while preserving immediate visual feedback and alignment with programming concepts. 
Such an approach is consistent with research recommendations calling for structured, accessible tools that enable learners to focus on logic and problem solving rather than technical implementation details \cite{md_shamsudin_crafting_2024}.
